\documentclass[conference]{IEEEtran}

\usepackage{cite}
\usepackage{amsmath}
\usepackage{booktabs}
\usepackage{graphicx}
\usepackage{url}

\newcommand{\pipesense}{PipeSense-ARM}

\begin{document}

\title{\pipesense: Safe Runtime Pipeline Adaptation in an Educational Embedded Core}

\author{
\IEEEauthorblockN{Tyler Lee}
\IEEEauthorblockA{\textit{Electrical and Computer Engineering} \\
\textit{NYU Tandon School of Engineering}\\
Brooklyn, NY 11201, USA \\
tl2966@nyu.edu}
}

\maketitle

\begin{abstract}
Static pipeline policies cannot match every phase of an embedded workload.
\pipesense{} is a SystemVerilog research artifact that places a rolling-window
event observer, a hysteretic mode controller, and a drain-before-switch
reconfiguration unit around a five-stage ARM-like educational core. The
controller selects among normal, branch, memory, hazard, and low-activity
modes without benchmark labels or software intervention. Across 13
embedded-style workloads and six configurations, adaptive execution reduces
total cycles by 7.59\% and an activity-energy proxy by 5.78\% relative to
static normal mode; nine workloads improve, three tie, and one regresses. The
adaptive policy nevertheless remains an average 8.90\% behind the hindsight
best fixed mode, exposing policy headroom. All 78 directed runs match a
sequential ISA reference and report zero safety faults or timeouts. A 500-seed
constrained-random regression produces 2,500 mode-result rows with zero
assertion failures, safety faults, or timeouts. Yosys reports a 165.03\%
integrated generic-cell overhead proxy, showing that the initial observer is
not yet cost-effective for a tiny core. The result is a reproducible control
artifact with explicit performance, safety, and cost limits rather than an
ARM-compatible processor or calibrated hardware implementation.
\end{abstract}

\begin{IEEEkeywords}
adaptive microarchitecture, embedded processors, hardware monitoring,
pipeline reconfiguration, SystemVerilog
\end{IEEEkeywords}

\section{Introduction}

Embedded programs can move between control-heavy loops, streaming memory
accesses, dependency chains, and low-utilization intervals. Program phases
have long motivated sampling and reconfigurable hardware
\cite{sherwood2002phase,dhodapkar2002working}. A fixed pipeline policy is easy
to verify, but it can retain a branch penalty during a branch-heavy interval or
pay memory waits when a local mitigation would be more useful. Software and
firmware can react to coarse counters, but a small in-core controller can
observe the causes of stalls directly and respond without an interrupt or
runtime service.

This paper studies a narrow question: can a hardware-resident feedback loop
classify pipeline behavior, select a matching microarchitectural mode, and
change modes without corrupting in-flight execution? \pipesense{} implements
that loop in a five-stage IF/ID/EX/MEM/WB processor model. The observer counts
branches, memory waits, load-use hazards, stalls, valid stages, and retirement
over short windows. A controller applies stability and residency conditions,
and a separate reconfiguration unit gates fetch and drains the pipeline before
committing a requested mode.

The contribution is the integration and evaluation of this inspectable loop,
not a claim that its individual mechanisms are new. The artifact contributes:
1) an event-driven phase observer and hysteretic controller attached to a
sequential RTL core; 2) an explicit drain-before-switch contract checked by
simulation assertions, architectural reference comparison, and randomized
stress; and 3) an evaluation that preserves unfavorable results through fixed
baselines, a hindsight oracle, parameter sweeps, ablations, and synthesis cost
proxies. This combination exposes both the benefit and the present cost of
hardware-native adaptation.

\section{Background and Related Work}

Branch prediction and memory buffering are established ways to reduce control
and memory penalties \cite{smith1981branch,jouppi1990victim}. Reconfigurable
memory hierarchies further show that a hardware structure can trade
performance and energy as behavior changes
\cite{balasubramonian2000memory}. Phase-characterization work uses recurring
program signatures to select representative behavior or manage multiple
configurations \cite{sherwood2002phase,dhodapkar2002working}. \pipesense{}
uses the same broad observation that behavior changes over time, but it does
not construct basic-block or working-set signatures. It classifies direct
pipeline events with threshold logic over a short window.

Architectural simulators and activity models support early design exploration
\cite{austin2002simplescalar,brooks2000wattch}. Their results still depend on
the realism of the modeled machine and workload. Accordingly, the energy value
reported here is a sum of active pipeline slots with a fixed activity term
outside low-activity mode; it is not power or energy measured from hardware.
Likewise, the branch and memory modes are experimental mechanisms rather than
a production predictor or cache.

Runtime reconfiguration adds a correctness obligation beyond performance.
SystemVerilog assertions provide clocked checks close to RTL behavior
\cite{ieee1800_2023,foster2004assertion}. \pipesense{} combines such checks
with instruction tags and a sequential ISA model. The novelty boundary is
therefore modest: a compact artifact connects direct phase observation,
hysteretic control, conservative switching, and reproducible negative as well
as positive evidence in one educational pipeline.

\section{Design and Methodology}

\subsection{Core and Adaptive Modes}

The modeled core supports 32-bit educational encodings for arithmetic and
logical operations, load, store, compare, conditional branch, no-operation,
and a simulation-only halt sentinel. It includes valid bits, forwarding,
load-use interlocking, branch flushes, deterministic memory waits, writeback,
and performance counters. It is ARM-like in organization and instruction style
but is not ARM compatible.

Five modes expose distinct timing or activity behavior. Normal mode uses the
default branch, memory, and hazard paths. Branch mode resolves eligible
branches in decode and reduces the modeled flush from two slots to one. Memory
mode mitigates periodic data-memory wait requests, representing a small buffer
or prefetch effect. Hazard mode enables an additional load-result bypass and
suppresses the corresponding load-use stall. Low-activity mode removes a fixed
term from the activity proxy. These modes make controller decisions observable
while keeping the prototype small; they should not be interpreted as complete
commercial mechanisms.

\begin{figure}[t]
\centering
\fbox{\begin{minipage}{0.94\columnwidth}
\centering
\textbf{IF} $\rightarrow$ \textbf{ID} $\rightarrow$ \textbf{EX}
$\rightarrow$ \textbf{MEM} $\rightarrow$ \textbf{WB}\\[2pt]
$\downarrow$ narrow event taps\\[2pt]
\texttt{pipeline\_observer} $\rightarrow$
\texttt{adaptive\_controller}\\[-1pt]
$\downarrow$ mode request\\[-1pt]
\texttt{reconfig\_unit} $\rightarrow$ current mode
\end{minipage}}
\caption{Closed-loop organization of \pipesense{}. The policy observes only
pipeline events; mode changes are committed by a separate safety unit.}
\label{fig:architecture}
\end{figure}

\subsection{Observer and Controller}

Figure~\ref{fig:architecture} shows the feedback path. The observer samples
stage occupancy, branch-taken and flush events, load-use hazards, memory waits,
front-end stalls, cycles, and retired instructions. At the end of each
parameterized observation window, counters are compared with thresholds for
branch-heavy, memory-stall, load-use-hazard, front-end-stall, and
idle/low-utilization behavior. Otherwise, the window is labeled balanced. The
observer resets its counters after classification, so decisions depend on
recent events rather than a cumulative program profile.

The controller maps branch-heavy and front-end-stall phases to branch mode,
memory-stall phases to memory mode, load-use phases to hazard mode, and
idle/low-utilization phases to low-activity mode. Balanced behavior is sticky:
it retains the current mode instead of forcing a switch to normal. A stable
phase counter prevents one anomalous window from immediately changing policy.
Mode-specific minimum residency requirements provide additional hysteresis;
memory mode can be entered sooner than modes intended for less dominant events.
The observer window, event thresholds, and minimum residency are parameters
used in the sensitivity study.

\subsection{Safe Reconfiguration}

A controller request does not directly update the visible mode. The
reconfiguration unit first stops new fetches. Existing instructions continue
through IF/ID, ID/EX, EX/MEM, and MEM/WB until all four pipeline registers are
invalid and no memory wait is active. The unit then commits the latched mode,
acknowledges completion, and records both the event and its stall cycles. The
safe-boundary predicate is
\begin{equation}
\mathit{safe}=\mathit{pipeline\_empty}\land\neg\mathit{mem\_wait}.
\end{equation}
This conservative protocol can cost more than selective quiescence, but it
prevents an instruction from crossing a mode change that alters branch,
forwarding, or memory timing.

Simulation monitors check that retirement tags increase monotonically, no two
valid stages contain the same dynamic tag, mode changes coincide with the safe
boundary and completion acknowledgement, fetch remains gated during a drain,
the mode remains stable while reconfiguration is active, and the stall remains
bounded. The HDL result for each workload and mode is also compared with a
sequential ISA model using retired count and final architectural-state hash.
Formal harnesses exist for the isolated protocol, but no full-core formal proof
is claimed.

\section{Evaluation}

The directed matrix contains 13 workloads: arithmetic-heavy, branch-heavy,
memory-heavy, load-use-heavy, mixed-control, tiny-FIR, Dhrystone-style and
CoreMark-style toy ports, generated-style FIR and PID loops, long-FIR stress,
PID-phase stress, and randomized-memory-latency stress. The original
microbenchmarks isolate specific causes; the longer programs reduce dependence
on tiny loops without pretending to be full benchmark-suite ports. Every
workload runs in static normal, fixed branch, fixed memory, fixed hazard, fixed
low-activity, and adaptive configurations, producing 78 directed rows.

Metrics are cycles, retired instructions, instructions per cycle (IPC), stall
and flush cycles, memory waits, load-use stalls, reconfiguration events and
penalty, activity-energy proxy, safety faults, timeout status, and
architectural hash. Adaptive mode is compared with static normal and with the
best fixed mode chosen after observing each workload. The latter is not a
deployable policy; it is a hindsight oracle that reveals remaining opportunity.

The parameter study crosses observation windows of 16, 32, and 64 cycles;
minimum residencies of 8, 24, and 48 cycles; and tight, medium, and loose
threshold profiles. The resulting 27 configurations each rerun all 78 rows,
for 2,106 simulations. Ablations disable the observer, disable the controller,
or analytically remove measured reconfiguration penalty. A separate generator
runs legal constrained-random programs for 500 seeds and five execution modes.
The Docker evaluation uses Icarus Verilog 12.0 and Yosys 0.33. Scripts validate
row coverage, timeouts, safety fields, reference-model equality, sweep output,
and the public result summary. Source and commands are available at
\url{https://github.com/Tylerlee102/PipeSense-ARM}.

\section{Results}

\subsection{Performance and Oracle Comparison}

All 78 directed runs terminate, report zero safety faults, and pass the
retired-count and architectural-hash comparison. Table~\ref{tab:performance}
reports the complete cycle result. Adaptive execution uses 3,056 cycles versus
3,307 for static normal, a 7.59\% reduction. Average per-workload cycle
reduction is 6.33\%, average IPC improvement is 7.71\%, and the total
activity-energy proxy falls from 19,162 to 18,055, or 5.78\%. Nine workloads
improve, three tie, and generated DSP-FIR regresses by 1.04\% because two mode
changes cost eight cycles in a short run.

\begin{table}[t]
\centering
\caption{Adaptive cycle results. Oracle gap is relative to the best fixed mode;
the final oracle entry is the mean of per-workload gaps.}
\label{tab:performance}
\footnotesize
\begin{tabular}{lrrrr}
\toprule
Workload & Normal & Adapt. & Red. (\%) & Oracle (\%) \\
\midrule
Arithmetic       & 30  & 30  & 0.00  & -3.45 \\
Branch           & 128 & 114 & 10.94 & -8.57 \\
CoreMark-toy     & 162 & 162 & 0.00  & -8.00 \\
Dhrystone-toy    & 133 & 133 & 0.00  & -7.26 \\
DSP-FIR          & 192 & 194 & -1.04 & -15.48 \\
Load-use         & 205 & 178 & 13.17 & -5.33 \\
Long-FIR         & 876 & 792 & 9.59  & -1.80 \\
Memory           & 231 & 162 & 29.87 & -10.20 \\
Mixed-control    & 131 & 127 & 3.05  & -20.95 \\
PID-codegen      & 192 & 188 & 2.08  & -5.62 \\
PID-phase        & 653 & 620 & 5.05  & -1.47 \\
Random-memory    & 215 & 204 & 5.12  & -7.94 \\
Tiny-FIR         & 159 & 152 & 4.40  & -19.69 \\
\midrule
Total / mean     & 3307 & 3056 & 7.59 & -8.90 \\
\bottomrule
\end{tabular}
\end{table}

The oracle result limits the performance claim. Adaptive execution is behind
the best fixed choice on every workload, with a mean reported gap of 8.90\%.
Long-FIR and PID-phase come within 1.80\% and 1.47\%, respectively, while
mixed-control and tiny-FIR remain 20.95\% and 19.69\% behind. The controller
therefore captures useful phase behavior but is not an optimal online policy.
The large static-normal gain on memory-heavy (29.87\%) coexists with a 10.20\%
oracle gap, illustrating why both baselines are necessary.

\subsection{Sensitivity, Safety, and Cost}

The full sweep contains 1,755 adaptive-versus-fixed comparison cells. Adaptive
wins 478 and does not win 1,277. Against static normal alone, mean cycle
reduction increases from 1.91\% at a 16-cycle window to 3.72\% at 32 cycles and
4.29\% at 64 cycles, averaged across the other sweep dimensions. Tight,
medium, and loose thresholds average 3.98\%, 3.22\%, and 2.72\%, respectively.
These results show policy sensitivity rather than a universally robust default.

Table~\ref{tab:additional} summarizes the remaining evidence. Disabling either
the observer or controller returns adaptive execution to the 3,307-cycle
static-normal aggregate, so the improvement is not caused by a hidden default
mode. Removing all measured reconfiguration penalty would reduce cycles only
from 3,056 to 3,009, a further 1.54\%; phase selection is therefore a larger
opportunity than drain latency alone.

\begin{table}[t]
\centering
\caption{Ablation, safety, and generic synthesis evidence.}
\label{tab:additional}
\footnotesize
\begin{tabular}{lrr}
\toprule
Evidence & Result & Relative outcome \\
\midrule
Observer disabled & 3307 cycles & +8.21\% cycles \\
Controller disabled & 3307 cycles & +8.21\% cycles \\
Zero-cost switching & 3009 cycles & -1.54\% cycles \\
Safety fuzz & 2500 rows & 0 failures/timeouts \\
Baseline core proxy & 1830 cells & reference \\
Integrated proxy & 4850 cells & +165.03\% cells \\
\bottomrule
\end{tabular}
\end{table}

The 500-seed fuzz run reports zero assertion failures, safety faults, timeouts,
or nonzero simulator exits across 2,500 mode-result rows. This is stress
evidence, not exhaustive verification. Yosys maps the baseline proxy to 1,830
generic cells and the integrated proxy to 4,850, a 165.03\% delta. Standalone
observer, controller, and reconfiguration modules total 2,885 cells; the
observer alone accounts for 2,128. These figures are neither FPGA utilization
nor ASIC area, but they identify observer cost as a central design weakness.

\section{Discussion and Limitations}

The results support a scoped conclusion: direct pipeline events can drive a
closed-loop RTL controller that improves over a static normal policy while
preserving the modeled architectural and switching invariants. The negative
oracle result and synthesis proxy are equally important. They show that the
current classifier leaves substantial mode-selection headroom and that wide
rolling counters are disproportionately expensive beside a small teaching
core.

Internal validity is limited by stylized mode effects. Memory optimization
suppresses a deterministic wait request rather than modeling a cache or
prefetch queue. Branch optimization changes resolution timing rather than
implementing a predictor. The low-activity result is based on an activity sum,
not voltage, capacitance, frequency, or measured power. Yosys generic cells do
not establish timing closure or physical area. These proxies are suitable for
testing the control loop, not for claiming production speed, energy, or cost.

External validity is limited by the educational ISA and workload suite. The
toy and generated kernels are transparent and reference-checkable, but they do
not replace compiler-generated Embench, CoreMark, or application traces with a
real cache hierarchy, interrupts, and exceptions. Safety evidence is also
bounded: simulation assertions and randomized programs cover observed traces,
while the formal files prove only abstract or isolated properties and were not
rerun in the reported Docker environment. Full-core token conservation and
liveness remain future proof obligations.

A stronger implementation should replace counter banks with shared or sampled
counters, implement concrete branch and memory structures, tune policy on
training workloads separate from evaluation workloads, and measure area,
timing, and power on a named FPGA or standard-cell target. The existing oracle,
sweep, ablation, reference, and safety checks provide a basis for judging those
changes without discarding unfavorable rows.

\section{Conclusion}

\pipesense{} demonstrates a complete hardware-resident adaptation loop in a
small sequential processor model. A rolling-window observer identifies local
bottlenecks, a hysteretic controller requests a matching mode, and a separate
unit commits changes only after a safe drain. The adaptive configuration
reduces total cycles by 7.59\% and the activity-energy proxy by 5.78\% relative
to static normal while passing directed reference checks and randomized safety
stress. Its 8.90\% mean gap to the hindsight best fixed mode and 165.03\%
generic-cell overhead prevent a broader efficiency claim. The artifact is
therefore most useful as a reproducible foundation for safer and more
cost-effective adaptive embedded microarchitectures.

\bibliographystyle{IEEEtran}
\IEEEtriggeratref{2}
\bibliography{references}

\end{document}
